\begin{figure*}[t]
\providecommand{\width}{}
\renewcommand{\width}{0.28\textwidth}
\providecommand{\path}{}
\renewcommand{\path}{score_agent}
\captionsetup[subfigure]{position=above,labelformat=empty,margin={1em,0em},skip=-.1em}
\centering\hfil\hfil%
\begin{subfigure}[t]{\width}
\centering
\caption{Cartpole Swing Up}
\includegraphics[width=\textwidth]{\path/cartpole_swingup}
\end{subfigure}\hfil%
\begin{subfigure}[t]{\width}
\centering
\caption{Reacher Easy}
\includegraphics[width=\textwidth]{\path/reacher_easy}
\end{subfigure}\hfil%
\begin{subfigure}[t]{\width}
\centering
\caption{Cheetah Run}
\includegraphics[width=\textwidth]{\path/cheetah_run}
\end{subfigure}\hfil\hfil\hfil \\
\vspace{.7em}\hfil\hfil%
\begin{subfigure}[t]{\width}
\centering
\caption{Finger Spin}
\includegraphics[width=\textwidth]{\path/finger_spin}
\end{subfigure}\hfil%
\begin{subfigure}[t]{\width}
\centering
\caption{Cup Catch}
\includegraphics[width=\textwidth]{\path/cup_catch}
\end{subfigure}\hfil%
\begin{subfigure}[t]{\width}
\centering
\caption{Walker Walk}
\includegraphics[width=\textwidth]{\path/walker_walk}
\end{subfigure}\hfil\hfil\hfil \\
\vspace{1em}\hspace{1em}
\begin{subfigure}[t]{.7\textwidth}
\centering
\includegraphics[width=\textwidth]{\path/legend}
\end{subfigure}\hfil%
\caption{不同智能体设计的比较。图中展示测试性能随已收集回合数的变化。我们比较 \method、一个使用随机动作收集数据的版本（random collection），以及一个在每个环境步从 $1000$ 条动作序列中选择最优序列、但不通过 CEM 迭代细化规划的版本（random shooting）。曲线表示中位数，阴影区域表示在 5 个随机种子和 10 条轨迹上得到的第 5 到第 95 百分位范围。}
\label{fig:score_agent}
\end{figure*}
